\subsection*{4.4.}
回顾 $\mathbb Q_p(\xi_p)$ 的整数环的剩余域为 $\mathbb F_p$. 设

\begin{equation}\tag{4.4.0}
K=\text{$\mathbb Q_p(\xi_p)$ 的剩余域为 $\mathbb F_q$ 的非分歧扩张},
\end{equation}

并设

\begin{equation}\tag{4.4.1}
\tau\in\operatorname{Gal}(K/\mathbb Q_p(\xi_p))
\end{equation}

为“Frobenius”, 即约化后为 $\mathbb F_q/\mathbb F_p$ 的 Frobenius 自同构
$x\longmapsto x^p$ 的自同构.

设

\begin{equation}\tag{4.4.2}
F(X)=\sum_w A_wX^w\in K[X_1,\ldots,X_{n+2}]
\end{equation}

为 $f(X)$ 的提升, 其系数满足 $A_w^q=A_w$.
结合 (4.3.1) 及 (4.0.9), 并考虑 $\chi$ 的定义 (4.1.0), 得到

\begin{equation}\tag{4.4.3}
\begin{aligned}
&(q^\nu-1)\operatorname{card}V^*(\mathbb F_{q^\nu})
-\frac{(q^\nu-1)^{n+2}}{q^\nu}\\
&\quad=\frac1{q^\nu}
\sum_{(X_0,X)\in(\mu_{q^\nu-1}(\Omega))^{n+3}}
\prod_w\prod_{j=0}^{\operatorname{ord}(q^\nu)-1}
\Theta\!\left(A_w^{p^j}X_0^{p^j}X^{p^jw}\right).
\end{aligned}
\end{equation}

定义

\[
H(X)=H(X_0,\ldots,X_{n+2})\in K[[X_0,\ldots,X_{n+2}]]
\]

其定义式为

\begin{equation}\tag{4.4.4}
H(X)=\prod_w
\Theta\!\left(A_wX_0X_1^{w_1}\cdots X_{n+2}^{w_{n+2}}\right).
\end{equation}

令 $\operatorname{Gal}(K/\mathbb Q_p(\xi_p))$ 通过在系数上的作用
作用于 $K[[X_0,\ldots,X_{n+2}]]$;
$\sigma\in\operatorname{Gal}(K/\mathbb Q_p(\xi_p))$ 的逐分量作用记为

\begin{equation}\tag{4.4.5}
G(X)\longmapsto G^\sigma(X).
\end{equation}

特别地,

\[
H^{\tau^j}(X)=\prod_w
\Theta\!\left(A_w^{\tau^j}X_0X_1^{w_1}\cdots
X_{n+2}^{w_{n+2}}\right).
\]

由于 $A_w^{\tau^j}=A_w^{p^j}$, 公式 (4.4.3) 因此可改写为

\begin{equation}\tag{4.4.6}
\begin{aligned}
(q^\nu-1)\operatorname{card}V^*(\mathbb F_{q^\nu})
&=\frac{(q^\nu-1)^{n+2}}{q^\nu}\\
&\quad+\frac1{q^\nu}
\sum_{(X_0,X)\in(\mu_{q^\nu-1}(\Omega))^{n+3}}
\prod_{j=0}^{\operatorname{ord}(q^\nu)-1}H^{\tau^j}(X^{p^j}).
\end{aligned}
\end{equation}
